\textbf{Phần 1. Trắc nghiệm lựa chọn}
\setcounterpageref{ex}{0}

\begin{ex} %[2D1N06]
Số gồm $7$ chục và $0$ đơn vị là
\choice
{$7$}
{$77$}
{$17$}
{\True $70$}
\loigiai{
Số đó là $70$.\\
\textbf{Đáp án của câu hỏi là D}
}
\end{ex}

\begin{ex} %[2D1N06]
Số liền sau của số $89$ là
\choice
{$88$}
{$91$}
{\True $90$}
{$99$}
\loigiai{
$89 + 1 = 90$.\\
\textbf{Đáp án của câu hỏi là C}
}
\end{ex}

\begin{ex} %[2D1H06]
Trong phép trừ $75 - 32 = 43$, số trừ là
\choice
{$75$}
{$43$}
{$107$}
{\True $32$}
\loigiai{
Số trừ là 32.\\
\textbf{Đáp án của câu hỏi là D}
}
\end{ex}

\begin{ex} %[2D1H06]
Kết quả của phép tính $43 + 36$ là
\choice
{\True $79$}
{$78$}
{$80$}
{$69$}
\loigiai{
$43 + 36 = 79$.\\
\textbf{Đáp án của câu hỏi là A}
}
\end{ex}

\begin{ex} %[2D1H06]
Kết quả của phép tính $96 - 54$ là
\choice
{$41$}
{$43$}
{\True $42$}
{$44$}
\loigiai{
$96 - 54 = 42$.\\
\textbf{Đáp án của câu hỏi là C}
}
\end{ex}

\begin{ex} %[2D1H06]
Số 58 hơn số 20 bao nhiêu đơn vị?
\choice
{$28$}
{\True $38$}
{$48$}
{$78$}
\loigiai{
$58 - 20 = 38$.\\
\textbf{Đáp án của câu hỏi là B}
}
\end{ex}

\begin{ex} %[2D1V06]
Sắp xếp các số $42, 24, 65, 56$ theo thứ tự từ lớn đến bé:
\choice
{\True $65, 56, 42, 24$}
{$24, 42, 56, 65$}
{$65, 42, 56, 24$}
{$56, 65, 42, 24$}
\loigiai{
Thứ tự giảm dần: $65 > 56 > 42 > 24$.\\
\textbf{Đáp án của câu hỏi là A}
}
\end{ex}

\begin{ex} %[2D1V06]
Một sợi dây dài $65\text{ cm}$, người ta cắt bớt đi $23\text{ cm}$. Sợi dây còn lại dài là
\choice
{$41\text{ cm}$}
{\True $42\text{ cm}$}
{$43\text{ cm}$}
{$44\text{ cm}$}
\loigiai{
$65 - 23 = 42\text{ cm}$.\\
\textbf{Đáp án của câu hỏi là B}
}
\end{ex}

\begin{ex} %[2D1V06]
Tính giá trị biểu thức: $32 + 14 + 20$.
\choice
{\True $66$}
{$64$}
{$68$}
{$70$}
\loigiai{
$46 + 20 = 66$.\\
\textbf{Đáp án của câu hỏi là A}
}
\end{ex}

\begin{ex} %[2D1V06]
Tìm $x$ biết: $x - 30 = 25 + 14$.
\choice
{$68$}
{$70$}
{$59$}
{\True $69$}
\loigiai{
$x - 30 = 39 \Rightarrow x = 39 + 30 = 69$.\\
\textbf{Đáp án của câu hỏi là D}
}
\end{ex}

\begin{ex} %[2D1V06]
Lớp 2A có 20 bạn nam và 15 bạn nữ. Hỏi lớp 2A có tất cả bao nhiêu học sinh?
\choice
{$34$ bạn}
{\True $35$ bạn}
{$36$ bạn}
{$37$ bạn}
\loigiai{
$20 + 15 = 35$ bạn.\\
\textbf{Đáp án của câu hỏi là B}
}
\end{ex}

\begin{ex} %[2D1C06]
Tìm số có hai chữ số biết chữ số hàng chục gấp đôi chữ số hàng đơn vị và tổng hai chữ số bằng 9.
\choice
{$36$}
{$42$}
{\True $63$}
{$84$}
\loigiai{
$6 = 3 \times 2$ và $6 + 3 = 9$. Số đó là $63$.\\
\textbf{Đáp án của câu hỏi là C}
}
\end{ex}

\textbf{Phần 2. Trắc nghiệm đúng sai}
\setcounterpageref{ex}{0}

\begin{ex} %[2D1H06]
Xét tính đúng sai về số học và phép tính:
\choiceTF
{\True Số liền trước của số 100 là 99.}
{\True $40 + 50 = 90$.}
{Số 75 gồm 5 chục và 7 đơn vị.}
{\True Số 75 gồm 7 chục và 5 đơn vị.}
\loigiai{
\begin{itemize}
    \item $100 - 1 = 99$: Đúng.
    \item $40 + 50 = 90$: Đúng.
    \item 75 gồm 7 chục và 5 đơn vị (không phải 5 chục). Ý 3 Sai, Ý 4 Đúng.
\end{itemize}
\begin{itemize}
    \item a) \textbf{Đúng}.
    \item b) \textbf{Đúng}.
    \item c) \textbf{Sai}.
    \item d) \textbf{Đúng}.
\end{itemize}
}
\end{ex}

\begin{ex} %[2D1V06]
Một trang trại có 46 con thỏ trắng và 32 con thỏ nâu. Xét tính đúng sai:
\choiceTF
{\True Trang trại có tất cả 78 con thỏ ($46 + 32 = 78$).}
{\True Số thỏ trắng nhiều hơn thỏ nâu là 14 con ($46 - 32 = 14$).}
{Số thỏ nâu nhiều hơn thỏ trắng 14 con.}
{\True Thỏ trắng nhiều hơn thỏ nâu nên thỏ nâu ít hơn thỏ trắng 14 con.}
\loigiai{
\begin{itemize}
    \item Tổng: $46 + 32 = 78$: Đúng.
    \item Trắng hơn nâu: $46 - 32 = 14$: Đúng.
    \item Thỏ nâu ít hơn (không phải nhiều hơn). Ý 3 Sai, Ý 4 Đúng.
\end{itemize}
\begin{itemize}
    \item a) \textbf{Đúng}.
    \item b) \textbf{Đúng}.
    \item c) \textbf{Sai}.
    \item d) \textbf{Đúng}.
\end{itemize}
}
\end{ex}

\textbf{Phần 3. Trả lời ngắn}
\setcounterpageref{ex}{0}

\begin{ex} %[2D1H06]
Viết số liền trước của số $60$.
\shortans{59}
\loigiai{
\textbf{Đáp số:} $59$\\
$60 - 1 = 59$.
}
\end{ex}

\begin{ex} %[2D1H06]
Tính tổng: $53 + 24$.
\shortans{77}
\loigiai{
\textbf{Đáp số:} $77$\\
$53 + 24 = 77$.
}
\end{ex}

\begin{ex} %[2D1V06]
Tính hiệu: $78 - 35$.
\shortans{43}
\loigiai{
\textbf{Đáp số:} $43$\\
$78 - 35 = 43$.
}
\end{ex}

\begin{ex} %[2D1V06]
Tìm $x$ biết: $x - 24 = 43$.
\shortans{67}
\loigiai{
\textbf{Đáp số:} $67$\\
$x = 43 + 24 = 67$.
}
\end{ex}

\begin{ex} %[2D1V06]
Số lớn nhất có hai chữ số hơn số tròn chục lớn nhất có hai chữ số bao nhiêu đơn vị?
\shortans{9}
\loigiai{
\textbf{Đáp số:} $9$\\
$99 - 90 = 9$.
}
\end{ex}

\begin{ex} %[2D1C06]
Tìm số tròn chục nhỏ nhất có hai chữ số mà tổng hai chữ số bằng 5.
\shortans{50}
\loigiai{
\textbf{Đáp số:} $50$\\
Số tròn chục có dạng $\overline{a0} \Rightarrow a = 5 \Rightarrow 50$.
}
\end{ex}
